\chapter{The Edit Command}
\label{ch:edit}

\index{Edit command}
\index{sam}
\index{structural regular expressions}

The \cmd{Edit} command gives acme a complete text-processing language,
inherited from Rob Pike's sam editor.  It is the most powerful feature
in acme and rewards study, though it is not required for daily use.

To use it, type \cmd{Edit} followed by a sam command in any tag or
body, and middle-click (\button{2}) the whole thing.  For example:

\begin{Verbatim}[frame=single]
Edit ,s/old/new/g
\end{Verbatim}

\noindent
This replaces all occurrences of ``old'' with ``new'' in the current
window.

\section{Addresses}
\label{sec:addresses}

\index{addresses}

An address identifies a substring of the file.  Addresses appear before
commands to specify what text the command operates on.

\begin{longtable}{@{}lp{3.5in}@{}}
\toprule
Address & Meaning \\
\midrule
\endhead
\texttt{.} & Dot: the current selection. \\
\texttt{0} & The empty string at the beginning of the file. \\
\texttt{\$} & The empty string at the end of the file. \\
\texttt{\textit{n}} & Line \textit{n}. \\
\texttt{\#\textit{n}} & Character position \textit{n} (counting from 0). \\
\texttt{/\textit{regexp}/} & Next match of \textit{regexp} after dot. \\
\texttt{?\textit{regexp}?} & Previous match of \textit{regexp} before dot. \\
\texttt{\textit{a},\textit{b}} & From address \textit{a} to address \textit{b}. \\
\texttt{\textit{a};\textit{b}} & From \textit{a} to \textit{b}, with dot set to \textit{a}
  before evaluating \textit{b}. \\
\texttt{\textit{a}+\textit{b}} & Address \textit{b} evaluated from the end of \textit{a}. \\
\texttt{\textit{a}-\textit{b}} & Address \textit{b} evaluated from the start of \textit{a},
  searching backward. \\
\bottomrule
\end{longtable}

\paragraph{Examples:}

\begin{itemize}
\item \texttt{1,5} --- lines 1 through 5.
\item \texttt{,} --- the entire file (shorthand for \texttt{0,\$}).
\item \texttt{.+1} --- the line after the current selection.
\item \texttt{/func/} --- the next occurrence of ``func''.
\item \texttt{0,/\^{}package/} --- from start of file to the line
      matching ``package''.
\end{itemize}

\section{Simple Commands}
\label{sec:simple}

\begin{longtable}{@{}lp{3.5in}@{}}
\toprule
Command & Action \\
\midrule
\endhead
\texttt{a/\textit{text}/} & \textbf{Append}: insert \textit{text}
  after the addressed text. \\
\texttt{i/\textit{text}/} & \textbf{Insert}: insert \textit{text}
  before the addressed text. \\
\texttt{c/\textit{text}/} & \textbf{Change}: replace the addressed
  text with \textit{text}. \\
\texttt{d} & \textbf{Delete}: delete the addressed text. \\
\texttt{s/\textit{regexp}/\textit{text}/} & \textbf{Substitute}:
  replace the first match of \textit{regexp} in the addressed text
  with \textit{text}. \\
\texttt{s/\textit{regexp}/\textit{text}/g} & Substitute all matches
  (global flag). \\
\texttt{m \textit{addr}} & \textbf{Move}: move the addressed text to
  after \textit{addr}. \\
\texttt{t \textit{addr}} & \textbf{Copy}: copy the addressed text to
  after \textit{addr}. \\
\texttt{p} & \textbf{Print}: display the addressed text (in
  \file{+Errors}). \\
\texttt{=} & Print the line number and character offset of the
  address. \\
\bottomrule
\end{longtable}

The delimiter for \texttt{a}, \texttt{i}, \texttt{c}, and \texttt{s}
can be any character, not just \texttt{/}.  Newlines within the text are
entered literally (the text continues on the next line, terminated by
another delimiter).

\paragraph{Examples:}

\begin{Verbatim}[frame=single]
Edit 1 i/// Copyright 2024\n/
Edit ,d
Edit ,s/colour/color/g
Edit 5 a/\nnew line after line 5\n/
Edit /ERROR/ p
\end{Verbatim}

\section{Structural Commands}
\label{sec:structural}

\index{structural commands}

The structural commands are what make sam's language truly powerful.
They apply commands to \emph{each match} of a pattern, rather than to
lines.

\begin{longtable}{@{}lp{3.5in}@{}}
\toprule
Command & Action \\
\midrule
\endhead
\texttt{x/\textit{regexp}/\textit{cmd}} & \textbf{Extract}: for each
  match of \textit{regexp} in the addressed text, set dot to the match
  and execute \textit{cmd}. \\
\texttt{y/\textit{regexp}/\textit{cmd}} & \textbf{Complement}: for
  each piece of text between matches of \textit{regexp}, set dot and
  execute \textit{cmd}. \\
\texttt{g/\textit{regexp}/\textit{cmd}} & \textbf{Guard}: execute
  \textit{cmd} only if the addressed text contains a match of
  \textit{regexp}. \\
\texttt{v/\textit{regexp}/\textit{cmd}} & \textbf{Inverse guard}:
  execute \textit{cmd} only if the addressed text does \emph{not}
  contain a match of \textit{regexp}. \\
\bottomrule
\end{longtable}

The key insight from Pike's ``Structural Regular Expressions'' paper is
that \texttt{x} replaces the traditional line-oriented approach.
Instead of ``for each line,'' you say ``for each match.''

\paragraph{Idiom: ``for each line'' is \texttt{x/.*\textbackslash n/}.}

Since \texttt{x} matches arbitrary patterns, it operates on structure
rather than lines.  You can loop over words, sentences, paragraphs,
function bodies, or any textual pattern.

\section{Grouping}
\label{sec:grouping}

\index{grouping commands}

Multiple commands can be grouped with braces:

\begin{Verbatim}[frame=single]
Edit ,x/.*\n/ {
    g/TODO/ p
}
\end{Verbatim}

\noindent
This prints every line containing ``TODO''.

\section{Multi-File Commands}
\label{sec:multifile}

\index{X command}
\index{Y command}

\begin{longtable}{@{}lp{3.5in}@{}}
\toprule
Command & Action \\
\midrule
\endhead
\texttt{X/\textit{regexp}/\textit{cmd}} & For each open file whose
  name matches \textit{regexp}, execute \textit{cmd}. \\
\texttt{Y/\textit{regexp}/\textit{cmd}} & For each open file whose
  name does \emph{not} match \textit{regexp}, execute \textit{cmd}. \\
\bottomrule
\end{longtable}

\noindent
These operate across all open windows, not just the current one.

\paragraph{Example:} Replace ``foo'' with ``bar'' in all open \file{.c}
files:

\begin{Verbatim}[frame=single]
Edit X/\.c$/ ,s/foo/bar/g
\end{Verbatim}

\section{Substitution Details}
\label{sec:subst}

\index{substitute command}

The substitute command supports backreferences:

\begin{itemize}
\item \texttt{\textbackslash 1}, \texttt{\textbackslash 2}, \ldots\
      refer to captured groups in the regexp.
\item \texttt{\&} refers to the entire match.
\end{itemize}

\paragraph{Example:} Swap the order of two-word variable names:

\begin{Verbatim}[frame=single]
Edit ,s/([a-z]+)_([a-z]+)/\2_\1/g
\end{Verbatim}

\section{Worked Examples}
\label{sec:examples}

These examples demonstrate common editing tasks using the Edit command.

\subsection{Delete All Blank Lines}

\begin{Verbatim}[frame=single]
Edit ,x/\n\n+/ c/\n/
\end{Verbatim}

\noindent
Translation: in the whole file (\texttt{,}), for each run of two or
more newlines (\texttt{x/\textbackslash n\textbackslash n+/}), replace
with a single newline (\texttt{c/\textbackslash n/}).

\subsection{Indent Selected Lines}

\begin{Verbatim}[frame=single]
Edit x/.*\n/ i/\t/
\end{Verbatim}

\noindent
For each line in the selection, insert a tab at the beginning.  (Select
the lines first, then run this Edit command.)

\subsection{Remove Trailing Whitespace}

\begin{Verbatim}[frame=single]
Edit ,s/[ \t]+$//g
\end{Verbatim}

\subsection{Rename a Variable}

\begin{Verbatim}[frame=single]
Edit ,s/\boldName\b/newName/g
\end{Verbatim}

\noindent
The \texttt{\textbackslash b} word boundaries prevent partial matches.

\subsection{Delete Lines Containing ``DEBUG''}

\begin{Verbatim}[frame=single]
Edit ,x/.*\n/ g/DEBUG/ d
\end{Verbatim}

\noindent
For each line, if it contains ``DEBUG'', delete it.

\subsection{Extract Function Signatures from a C File}

\begin{Verbatim}[frame=single]
Edit ,x/^[a-zA-Z].*\n{/ p
\end{Verbatim}

\noindent
Print every line that starts with a letter and ends with \texttt{\{},
which is a rough match for C function definitions.

\subsection{Number All Lines}

Select the text, then:

\begin{Verbatim}[frame=single]
Edit |awk '{printf "%4d  %s\n", NR, $0}'
\end{Verbatim}

\noindent
Sometimes the easiest ``Edit'' command is a pipe to a Unix tool.

\section{What Is Not Implemented}

Acme implements most of the sam command language, with these exceptions:

\begin{itemize}
\item \texttt{k} (mark) --- not implemented.
\item \texttt{n} (print file list) --- not implemented.
\item \texttt{q} (quit) --- use \cmd{Exit} instead.
\item \texttt{!} (shell command) --- use acme's native shell execution.
\end{itemize}

\section{Further Reading}

The Edit command language is fully described in:

\begin{itemize}
\item Rob Pike, ``The Text Editor sam,'' \emph{Software---Practice and
      Experience}, Vol.~17, No.~11, November 1987.
\item Rob Pike, ``Structural Regular Expressions,'' Proc.\ EUUG Spring
      1987 Conference.
\item The sam(1) manual page.
\end{itemize}
